% ch25_2categories_bicategories.tex — Volume III Chapter 1

\chapter{2-Categories and Bicategories}

\begin{overviewbox}
Chapter~22 introduced the 2-category $\mathbf{Cat}$ — objects are categories,
1-cells are functors, 2-cells are natural transformations — and proved the
interchange law for its vertical and horizontal compositions. This chapter
develops the full theory of 2-categorical structure: strict 2-categories,
\emph{bicategories} (where associativity and unitality hold only up to
coherent 2-isomorphism), the three flavours of functor between them
(\emph{strict}, \emph{pseudo}, \emph{lax}), and the higher cells
(\emph{pseudonatural transformations} and \emph{modifications}) that organize
them into a 3-categorical hierarchy.

Two payoffs. First, examples beyond $\mathbf{Cat}$ — $\mathbf{Span}$,
$\mathbf{Prof}$, $\mathbf{Rel}$, $\mathbf{Bimod}$ — none of which are
strictly associative; bicategory language is mandatory to handle them.
Second, the strictification theorem: every bicategory is biequivalent to a
strict 2-category. This is the higher-categorical analogue of Mac Lane's
coherence (Chapter~19) and explains why we may often pretend strictness is
no loss of generality.
\end{overviewbox}


% ═══════════════════════════════════════════════════════════════════════════
\section{Strict 2-Categories Revisited}
\label{sec:strict-2cat-revisited}
% ═══════════════════════════════════════════════════════════════════════════

\subsection{Formalizing}

\begin{definition}[Strict 2-category]
A \term{strict 2-category} $\mathcal{K}$ consists of:
\begin{itemize}
  \item A class of \emph{objects} (0-cells) $a, b, \ldots$
  \item For each pair $a, b$, a category $\mathcal{K}(a, b)$ — the
        \emph{hom-category}. Its objects are 1-cells $f : a \to b$; its
        morphisms are 2-cells $\alpha : f \Rightarrow g$.
  \item For each triple $a, b, c$, a \emph{composition functor}
        $\circ : \mathcal{K}(b,c) \times \mathcal{K}(a,b) \to \mathcal{K}(a,c)$.
  \item For each $a$, an identity 1-cell $1_a \in \mathcal{K}(a,a)$.
\end{itemize}
These satisfy \emph{strict} associativity and unit laws:
$h \circ (g \circ f) = (h \circ g) \circ f$ and $1_b \circ f = f = f \circ 1_a$.
\end{definition}

The composition functor being a \emph{functor} (not merely a function) is
what packages the interchange law: applied to 2-cells it gives horizontal
composition, applied to identity 2-cells it gives whiskering, and functoriality
in each variable says exactly that vertical and horizontal composites commute.

\begin{remark}[2-categories as $\mathbf{Cat}$-enriched categories]
A strict 2-category is precisely a category \emph{enriched over $\mathbf{Cat}$}
with the cartesian product. Hom-categories are the $\mathcal{V}$-objects;
composition functors are the $\mathcal{V}$-composition. All structure
descends from Chapter~12's enriched-category machinery.
\end{remark}

\subsection{Formally}

\begin{example_thm}
$\mathbf{Cat}$ itself, with hom-categories $\mathbf{Cat}(\mathcal{A}, \mathcal{B})
= [\mathcal{A}, \mathcal{B}]$ the functor category.
\end{example_thm}

\begin{example_thm}
For each ordinary category $\mathcal{C}$, the \emph{locally discrete} 2-category
$\mathcal{C}_{\mathrm{disc}}$: same objects and 1-cells; only identity 2-cells.
This realises every category as a 2-category.
\end{example_thm}

\begin{example_thm}
For each ordinary category $\mathcal{C}$ with all pullbacks, the 2-category
$\mathbf{Cat}/\mathcal{C}$ of \emph{categories sliced over $\mathcal{C}$}:
objects are functors into $\mathcal{C}$, 1-cells are triangles, 2-cells are
natural transformations between the triangle's tops compatible with the
projection.
\end{example_thm}

\begin{example_thm}[Monoidal categories as one-object 2-categories]
A monoidal category $(\mathcal{V}, \otimes, I)$ is the same data as a strict
2-category $B\mathcal{V}$ with a single object $\star$: 1-cells are objects
of $\mathcal{V}$, composition is $\otimes$, 2-cells are morphisms of
$\mathcal{V}$, the identity 1-cell is $I$. This works when $\otimes$ is
strictly associative; otherwise the identification is up to a coherent
isomorphism, which is the bicategorical case (\S\ref{sec:bicategories}).
\end{example_thm}


% ═══════════════════════════════════════════════════════════════════════════
\begin{nameorigin}
\term{Bicategory} was introduced by \emph{Jean Bénabou} in 1967. The
prefix \emph{bi-} (Latin ``two'') reflects that there are two levels of
morphisms (1-cells and 2-cells), but with associativity and unit laws
holding only ``up to coherent 2-isomorphism'' rather than strictly.
This is in contrast to the \emph{strict} 2-category, where the laws
hold on the nose. The strict-vs-weak distinction is one of the central
themes of higher category theory.

\term{Pseudofunctor} (a functor between bicategories that preserves
composition only up to coherent isomorphism) is also Bénabou's term;
the prefix \emph{pseudo-} (Greek ``false, almost'') indicates that the
preservation is up to coherent iso rather than strict equality.
\term{Lax functor}: even weaker, with only a (not necessarily
invertible) 2-cell as the comparison; ``lax'' from Latin
\emph{laxus} (``loose, slack'').
\end{nameorigin}

\section{Bicategories}
\label{sec:bicategories-formal}
% ═══════════════════════════════════════════════════════════════════════════

\subsection{Informally}

A bicategory is a 2-category in which the strict equality
$h \circ (g \circ f) = (h \circ g) \circ f$ is replaced by an
\emph{associator}: a specified invertible 2-cell
$\alpha_{h,g,f} : (h \circ g) \circ f \xRightarrow{\;\cong\;} h \circ (g \circ f)$.
Likewise the strict unit equations are replaced by invertible \emph{unitors}.
Coherence axioms — the pentagon for associators, the triangle for unitors —
guarantee that all reassociation diagrams commute, just as in the monoidal
case of Chapter~19.

\subsection{Formalizing}

\begin{definition}[Bicategory]
\label{def:bicategory}
A \term{bicategory} $\mathcal{B}$ consists of:
\begin{itemize}
  \item Objects (0-cells) $a, b, \ldots$
  \item For each pair $a, b$, a hom-category $\mathcal{B}(a,b)$.
  \item Composition functors $\circ_{a,b,c} : \mathcal{B}(b,c) \times \mathcal{B}(a,b)
        \to \mathcal{B}(a,c)$.
  \item Identity 1-cells $1_a \in \mathcal{B}(a,a)$.
  \item For each composable triple $(h, g, f)$, an invertible \term{associator}
        $\alpha_{h,g,f} : (h \circ g) \circ f \xRightarrow{\;\cong\;}
         h \circ (g \circ f)$, natural in $h, g, f$.
  \item For each $f : a \to b$, invertible \term{unitors}
        $\lambda_f : 1_b \circ f \xRightarrow{\;\cong\;} f$ and
        $\rho_f : f \circ 1_a \xRightarrow{\;\cong\;} f$, natural in $f$.
\end{itemize}
These are required to satisfy the \term{pentagon} and \term{triangle} axioms:
\begin{itemize}
  \item Pentagon: for composable $k, h, g, f$, the two ways to reassociate
        $((k \circ h) \circ g) \circ f \to k \circ (h \circ (g \circ f))$
        via $\alpha$'s agree.
  \item Triangle: $\rho_g \circ_h \mathrm{id}_h = \mathrm{id}_g \circ_h \lambda_h$
        composed with $\alpha_{g,1_b,h}$ for $g : b \to c$, $h : a \to b$.
\end{itemize}
\end{definition}

\begin{howtosay}
\begin{itemize}
  \item $\alpha_{h,g,f}$ \sayit{``alpha sub h, g, f''} or \sayit{``the associator
        at h, g, f''}.
  \item $f \circ_h g$ \sayit{``f horizontal g''} — horizontal composition of
        2-cells, the action of the composition \emph{functor} on morphisms.
\end{itemize}
\end{howtosay}

A strict 2-category is the special case where all associators and unitors are
identity 2-cells. Conversely:

\begin{proposition}[Coherence for bicategories — preview]
\label{prop:bicat-coherence}
In any bicategory, every diagram built from associators and unitors commutes.
\end{proposition}

This is the analogue of Mac Lane's coherence theorem (Chapter~19) for the
monoidal case. The proof, via strictification, appears in
\S\ref{sec:strictification}.

\subsection{Reorganizing: Monoidal Categories as One-Object Bicategories}

\begin{proposition}
A monoidal category $(\mathcal{V}, \otimes, I, \alpha, \lambda, \rho)$ is the
same data as a bicategory with a single object.
\end{proposition}

\begin{proof}
A one-object bicategory $\mathcal{B}$ with object $\star$ has a single
hom-category $\mathcal{B}(\star, \star)$ whose objects (1-cells) and morphisms
(2-cells) are precisely the objects and morphisms of a category $\mathcal{V}$.
Composition $\circ : \mathcal{B}(\star,\star) \times \mathcal{B}(\star,\star)
\to \mathcal{B}(\star,\star)$ is a functor $\otimes : \mathcal{V} \times
\mathcal{V} \to \mathcal{V}$. The identity 1-cell is the unit $I$. The
associator and unitors are precisely $\alpha$, $\lambda$, $\rho$. The pentagon
and triangle axioms transfer identically. The correspondence is exact.
\end{proof}

\begin{remark}
This is the same kind of one-object delooping as ``monoid = one-object
category'' and ``group = one-object groupoid''. Each level of categorification
adds one dimension to the diagram and one axis of coherence.
\end{remark}


% ═══════════════════════════════════════════════════════════════════════════
\section{Pseudofunctors, Lax Functors, and Oplax Functors}
\label{sec:pseudofunctors}
% ═══════════════════════════════════════════════════════════════════════════

\subsection{Formalizing}

For ordinary categories, functoriality is an equation: $F(g \circ f) = F(g)
\circ F(f)$ and $F(1_a) = 1_{Fa}$. For bicategories, equality is too rigid:
when $\mathcal{B}$ has non-strict associators, we cannot expect a map of
bicategories to preserve composition on the nose. The right notion is
preservation \emph{up to a coherent 2-isomorphism}.

\begin{definition}[Pseudofunctor]
\label{def:pseudofunctor}
A \term{pseudofunctor} (also: \term{weak 2-functor}) $F : \mathcal{B} \to
\mathcal{C}$ between bicategories consists of:
\begin{itemize}
  \item A function $F : \mathrm{ob}\,\mathcal{B} \to \mathrm{ob}\,\mathcal{C}$
        on objects.
  \item For each pair $a, b$, a functor $F_{a,b} : \mathcal{B}(a,b) \to
        \mathcal{C}(Fa, Fb)$.
  \item For each composable pair $(g, f)$, an invertible 2-cell
        $\phi_{g,f} : F(g) \circ F(f) \xRightarrow{\;\cong\;} F(g \circ f)$ —
        the \term{compositor}.
  \item For each $a$, an invertible 2-cell $\phi_a : 1_{Fa} \xRightarrow{\;\cong\;}
        F(1_a)$ — the \term{unitor}.
\end{itemize}
The compositor and unitor are required to be \emph{natural} in their arguments
and to satisfy three coherence diagrams: one matching $F(\alpha)$ to the
associator of $\mathcal{C}$ via $\phi$, and two matching $F(\lambda)$,
$F(\rho)$ to the unitors via $\phi$.
\end{definition}

\begin{definition}[Lax and oplax functors]
A \term{lax functor} $F : \mathcal{B} \to \mathcal{C}$ has the same data as a
pseudofunctor except $\phi_{g,f}$ and $\phi_a$ are \emph{not required to be
invertible} — they are merely 2-cells in the displayed direction. An
\term{oplax functor} reverses the direction: $\phi_{g,f}^{\mathrm{op}} :
F(g \circ f) \Rightarrow F(g) \circ F(f)$ and $\phi_a^{\mathrm{op}} : F(1_a)
\Rightarrow 1_{Fa}$.
\end{definition}

\begin{remark}
Lax functors with values in $\mathbf{Cat}$ have a striking interpretation
that returns in Chapter~32: a lax functor $\mathbf{1} \to \mathbf{Cat}$ is
precisely a \emph{monad}. The compositor $\phi$ is the multiplication
$\mu$; the unitor $\phi_a$ is $\eta$; the coherence diagrams are exactly
the monad associativity and unit axioms.
\end{remark}

\subsection{Reorganizing: When Pseudofunctors Are Strict}

A \term{strict 2-functor} is a pseudofunctor whose $\phi_{g,f}$ and $\phi_a$
are all identity 2-cells. Between strict 2-categories these are the ``honest''
2-functors. Between bicategories, asking for strictness is too much; even the
identity functor of a non-strict bicategory cannot be strict, because the
associator is non-identity. \emph{Pseudo} is the natural notion.

\subsection{Formally}

\begin{example_thm}[Identity pseudofunctor]
$\mathrm{Id}_\mathcal{B} : \mathcal{B} \to \mathcal{B}$ with $\phi_{g,f} =
\mathrm{id}$ and $\phi_a = \mathrm{id}$. This is strict — but only because
the source equals the target.
\end{example_thm}

\begin{example_thm}[Underlying ordinary category]
The forgetful pseudofunctor $U : \mathcal{B} \to \mathcal{B}_0$ to the
underlying ordinary category collapses 2-cells. It is \emph{lax}, with
$\phi_{g,f} : g \circ_{\mathcal{B}} f \to g \circ_{\mathcal{B}_0} f$ given by
quotienting; not invertible in general.
\end{example_thm}

\begin{example_thm}[Free monoidal functor]
The free symmetric monoidal completion $\mathcal{V} \to \mathrm{Sym}(\mathcal{V})$
is a pseudofunctor of monoidal-as-bicategories. Compositors arise from the
canonical isomorphisms in the symmetric monoidal closure.
\end{example_thm}


% ═══════════════════════════════════════════════════════════════════════════
\section{Pseudonatural Transformations and Modifications}
\label{sec:pseudonat-modifications}
% ═══════════════════════════════════════════════════════════════════════════

\subsection{Formalizing}

For ordinary functors $F, G : \mathcal{A} \to \mathcal{B}$, a natural
transformation $\eta : F \Rightarrow G$ is a family $\eta_a : Fa \to Ga$
satisfying a square equality. For pseudofunctors $F, G : \mathcal{B} \to
\mathcal{C}$ between bicategories, the square is replaced by an invertible
2-cell.

\begin{definition}[Pseudonatural transformation]
\label{def:pseudonat}
A \term{pseudonatural transformation} $\eta : F \Rightarrow G$ between
pseudofunctors $F, G : \mathcal{B} \to \mathcal{C}$ consists of:
\begin{itemize}
  \item For each $a \in \mathcal{B}$, a 1-cell $\eta_a : Fa \to Ga$ in
        $\mathcal{C}$.
  \item For each 1-cell $f : a \to b$ in $\mathcal{B}$, an invertible 2-cell
        \begin{equation*}
          \eta_f : Gf \circ \eta_a \xRightarrow{\;\cong\;} \eta_b \circ Ff.
        \end{equation*}
\end{itemize}
These satisfy two coherence axioms: \emph{respect for composition} ($\eta_{g
\circ f}$ assembles from $\eta_g$ and $\eta_f$ via the compositors of $F, G$)
and \emph{respect for identities} ($\eta_{1_a}$ matches the unitor data).
\end{definition}

\begin{definition}[Lax/oplax transformations]
Drop the invertibility of $\eta_f$ to get a \term{lax} transformation; reverse
its direction to get an \term{oplax} transformation.
\end{definition}

\subsection{Reorganizing}

\begin{definition}[Modification]
\label{def:modification}
For two pseudonatural transformations $\eta, \eta' : F \Rightarrow G$, a
\term{modification} $m : \eta \Rrightarrow \eta'$ is a family of 2-cells
$m_a : \eta_a \Rightarrow \eta'_a$ in $\mathcal{C}$ — one for each object
$a \in \mathcal{B}$ — satisfying compatibility with the structure 2-cells
$\eta_f, \eta'_f$: for each $f : a \to b$, the equation
\begin{equation*}
  (\mathrm{id}_{\eta'_b} \circ_h Ff) \cdot m_a^{Gf} \;=\; m_b^{Ff} \cdot
  (\mathrm{id}_{Gf} \circ_h \eta_a)
\end{equation*}
in $\mathcal{C}(Fa, Gb)$, with $m^{(-)}$ denoting whiskering of $m$ by
the indicated 1-cell.
\end{definition}

\begin{proposition}[The tricategorical hierarchy]
For any pair of bicategories $\mathcal{B}, \mathcal{C}$, pseudofunctors,
pseudonatural transformations, and modifications form a (strict) 2-category
$\mathrm{Bicat}(\mathcal{B}, \mathcal{C})$.
\end{proposition}

\begin{proof}[Sketch]
1-cell composition is pasting of 2-cells; 2-cell vertical and horizontal
compositions are inherited from $\mathcal{C}$. The interchange law and unit
laws of the strict 2-category structure follow from the corresponding axioms
in $\mathcal{C}$ together with the coherence axioms of pseudonaturality.
\end{proof}

\begin{remark}
Globally, bicategories, pseudofunctors, pseudonatural transformations, and
modifications form a \term{tricategory} $\mathbf{Bicat}$. The tricategorical
structure is needed because composition of pseudofunctors is associative only
up to a coherent pseudonatural transformation, not strictly.
\end{remark}


% ═══════════════════════════════════════════════════════════════════════════
\section{Examples Beyond $\mathbf{Cat}$}
\label{sec:bicat-examples}
% ═══════════════════════════════════════════════════════════════════════════

The strength of bicategory theory shows in examples where strict associativity
fails. We sketch four.

\subsection{$\mathbf{Span}(\mathcal{C})$}

For a category $\mathcal{C}$ with finite limits, the bicategory
$\mathbf{Span}(\mathcal{C})$ has:
\begin{itemize}
  \item Objects: objects of $\mathcal{C}$.
  \item 1-cells $a \to b$: spans $a \xleftarrow{p} x \xrightarrow{q} b$.
  \item 2-cells: morphisms of spans (a single map of apexes commuting with
        both legs).
  \item Composition: pullback. Given $a \leftarrow x \to b$ and $b \leftarrow
        y \to c$, the composite is $a \leftarrow x \times_b y \to c$.
\end{itemize}

\begin{remark}
Composition is associative only up to the canonical isomorphism between
iterated pullbacks $(x \times_b y) \times_c z \cong x \times_b (y \times_c z)$
— a genuine non-trivial associator. This is why $\mathbf{Span}$ cannot be
made strict (its strict version would require choosing pullbacks coherently,
which is rarely possible).
\end{remark}

\subsection{$\mathbf{Prof}$}

The bicategory of \term{profunctors}:
\begin{itemize}
  \item Objects: small categories.
  \item 1-cells $\mathcal{A} \to \mathcal{B}$: profunctors $P :
        \mathcal{B}^{\mathrm{op}} \times \mathcal{A} \to \mathbf{Set}$.
  \item 2-cells: natural transformations of profunctors.
  \item Composition: the coend formula $(Q \circ P)(c, a) = \int^{b}
        Q(c, b) \times P(b, a)$.
\end{itemize}

The associator for $\mathbf{Prof}$ encodes the Fubini theorem for coends
(Chapter~26). Composition is not strict because the iterated coends
$\int^b \int^c$ and $\int^c \int^b$ are isomorphic but distinct constructions.

\subsection{$\mathbf{Rel}$}

The bicategory of relations:
\begin{itemize}
  \item Objects: sets.
  \item 1-cells $A \to B$: relations $R \subseteq A \times B$.
  \item 2-cells: containment $R \subseteq R'$ (a poset structure).
  \item Composition: relational composition.
\end{itemize}

$\mathbf{Rel}$ is in fact \emph{locally posetal} — every hom-category is a
poset — and the composition is strictly associative on relations, so it can
be presented as a strict 2-category. We keep it under ``bicategories'' for
uniformity with $\mathbf{Span}$.

\subsection{$\mathbf{Bimod}$}

For commutative rings, the bicategory $\mathbf{Bimod}$:
\begin{itemize}
  \item Objects: commutative rings.
  \item 1-cells $R \to S$: $(S, R)$-bimodules.
  \item 2-cells: bimodule homomorphisms.
  \item Composition: tensor product $M \otimes_S N$.
\end{itemize}

The associator is the canonical isomorphism $(M \otimes_S N) \otimes_T P
\cong M \otimes_S (N \otimes_T P)$. Like $\mathbf{Span}$, this fails strict
associativity because the tensor product is only associative up to a chosen
isomorphism.


% ═══════════════════════════════════════════════════════════════════════════
\section{Biequivalences and Strictification}
\label{sec:strictification}
% ═══════════════════════════════════════════════════════════════════════════

\subsection{Formalizing}

\begin{definition}[Biequivalence]
A pseudofunctor $F : \mathcal{B} \to \mathcal{C}$ is a \term{biequivalence}
if it is:
\begin{itemize}
  \item \emph{Essentially surjective on objects}: every $c \in \mathcal{C}$
        is equivalent (via 1-cells with invertible 2-cells) to some $Fb$;
  \item \emph{Locally an equivalence}: each $F_{a,b} : \mathcal{B}(a,b) \to
        \mathcal{C}(Fa, Fb)$ is an equivalence of categories.
\end{itemize}
\end{definition}

\begin{remark}
This is the analogue of ``fully faithful and essentially surjective $\Rightarrow$
equivalence'' for ordinary functors, but at one categorical level higher. The
witness data is a pair $(F, G, \eta, \varepsilon)$ where $\eta, \varepsilon$
are pseudonatural \emph{equivalences} (themselves invertible up to invertible
modifications).
\end{remark}

\subsection{Formally}

\begin{theorem}[Strictification]
\label{thm:strictification}
Every bicategory $\mathcal{B}$ is biequivalent to a strict 2-category
$\mathcal{B}^{\mathrm{st}}$.
\end{theorem}

\begin{prooflog}
\proofstep{Define $\mathcal{B}^{\mathrm{st}}$ with the same objects as
$\mathcal{B}$; 1-cells $a \to b$ are finite \emph{strings} of composable
1-cells of $\mathcal{B}$ from $a$ to $b$ (including the empty string at $a$,
acting as identity).}{We replace composites by formal strings to gain strict
associativity.}
\proofstep{2-cells from one string to another are 2-cells in $\mathcal{B}$
between their respective composites $\widehat{(-)} : \mathrm{strings} \to
\mathrm{1\text{-}cells}$.}{The map $\widehat{(-)}$ takes a string $(f_1,
\ldots, f_n)$ to $f_n \circ \cdots \circ f_1$, left-associated.}
\proofstep{Composition of strings is concatenation; identity 1-cell at $a$
is the empty string.}{Both operations are strictly associative and unital
\emph{at the level of strings}.}
\proofstep{The pseudofunctor $\widehat{(-)} : \mathcal{B}^{\mathrm{st}} \to
\mathcal{B}$ sends a string to its left-associated composite.}{Provides
the candidate biequivalence.}
\proofstep{Compositor of $\widehat{(-)}$: for strings $s, t$, the canonical
$\widehat{t \cdot s} = \widehat{t} \circ \widehat{s}$ is built from
associators in $\mathcal{B}$. Invertible by definition.}{The associator chain
is well-defined by Mac Lane coherence (Proposition~\ref{prop:bicat-coherence}).}
\proofstep{$\widehat{(-)}$ is essentially surjective: each 1-cell $f \in
\mathcal{B}$ is the composite of the singleton string $(f)$.}{Trivially.}
\proofstep{$\widehat{(-)}$ is locally an equivalence: 2-cells between
strings, by construction, are 2-cells between their composites in
$\mathcal{B}$; so $\widehat{(-)}_{a,b}$ is fully faithful, and surjective on
objects by the singleton-string argument.}{Both fullness/faithfulness and
essential surjectivity at the hom-category level.}
\proofstep{Conclude: $\widehat{(-)} : \mathcal{B}^{\mathrm{st}} \to \mathcal{B}$
is a biequivalence.}{By definition of biequivalence.}
\end{prooflog}

\begin{corollary}[Mac Lane coherence for bicategories]
\label{cor:bicat-coherence}
Every diagram of associators and unitors in $\mathcal{B}$ commutes.
\end{corollary}

\begin{proof}
Such a diagram, transported via the strictification $\widehat{(-)}$, becomes
a diagram in the strict 2-category $\mathcal{B}^{\mathrm{st}}$ where
associators and unitors are identities. There, every diagram of identity
2-cells commutes trivially. Pulling back via $\widehat{(-)}$ (which is faithful
on 2-cells) gives commutativity in $\mathcal{B}$.
\end{proof}

\begin{remark}[Why strictification is not the whole story]
Strictification does not say every bicategory \emph{equals} a strict
2-category, only that it is biequivalent. The distinction matters when we
demand strict preservation of structure (as in strict $\Omega$-spectra of
algebraic topology). Mostly, however, the strictification theorem licenses
the working mathematician to write as if all bicategories were strict — the
``up to coherent isomorphism'' caveats can be left implicit.
\end{remark}


% ═══════════════════════════════════════════════════════════════════════════
\section{Exercises}
\label{sec:bicat-exercises}
% ═══════════════════════════════════════════════════════════════════════════

\begin{exercisesbox}
\begin{qa}
\qaitem{Define a strict 2-category in one line.}{A category enriched over $\mathbf{Cat}$ (with cartesian product).}
\qaitem{What is the hom-category of a strict 2-category $\mathcal{K}$?}{For each pair $a, b$, a category $\mathcal{K}(a,b)$ whose objects are 1-cells $a \to b$ and morphisms are 2-cells between them.}
\qaitem{What additional data does a bicategory have beyond a strict 2-category?}{Invertible 2-cells: an associator $\alpha_{h,g,f}$ for each composable triple, and left/right unitors $\lambda_f, \rho_f$ for each 1-cell.}
\qaitem{What axioms must the associator and unitors satisfy?}{The pentagon (for four-fold composites) and the triangle (for unitors).}
\qaitem{What is the bicategorical analogue of Mac Lane's coherence theorem?}{Every diagram of associators and unitors in a bicategory commutes (Corollary~\ref{cor:bicat-coherence}).}
\qaitem{What does a monoidal category correspond to as a bicategory?}{A one-object bicategory. The single hom-category is the underlying category of the monoidal category.}
\qaitem{State the strictification theorem.}{Every bicategory is biequivalent to a strict 2-category.}
\qaitem{What is a pseudofunctor?}{A map of bicategories preserving composition and identities up to invertible 2-cells (compositor $\phi_{g,f}$ and unitor $\phi_a$) satisfying coherence axioms.}
\qaitem{What is a lax functor?}{Same data as a pseudofunctor, but $\phi_{g,f}$ and $\phi_a$ are not required to be invertible.}
\qaitem{What is the difference between lax and oplax functors?}{Lax: $\phi_{g,f} : F(g) \circ F(f) \Rightarrow F(g \circ f)$. Oplax: arrows reversed.}
\qaitem{A lax functor $\mathbf{1} \to \mathbf{Cat}$ from the terminal bicategory is what?}{A monad. The compositor is $\mu$, the unitor is $\eta$, and the coherence axioms are the monad laws.}
\qaitem{What is a pseudonatural transformation?}{Between pseudofunctors $F, G$: a family of 1-cells $\eta_a : Fa \to Ga$ and an invertible 2-cell $\eta_f : Gf \circ \eta_a \cong \eta_b \circ Ff$ for each 1-cell $f$, satisfying coherence with composition and identities.}
\qaitem{What is a modification?}{Between pseudonatural transformations $\eta, \eta' : F \Rightarrow G$: a family of 2-cells $m_a : \eta_a \Rightarrow \eta'_a$ compatible with the structure 2-cells of $\eta$ and $\eta'$.}
\qaitem{Why are bicategories, pseudofunctors, pseudonatural transformations, and modifications a tricategory rather than a strict 3-category?}{Composition of pseudofunctors is associative only up to coherent pseudonatural transformation, not strictly.}
\qaitem{What is $\mathbf{Span}(\mathcal{C})$?}{The bicategory whose objects are objects of $\mathcal{C}$, 1-cells are spans $a \leftarrow x \to b$, 2-cells are span morphisms, and composition is pullback.}
\qaitem{Why is $\mathbf{Span}(\mathcal{C})$ not strictly associative?}{Pullbacks are only associative up to canonical isomorphism: $(x \times_b y) \times_c z \cong x \times_b (y \times_c z)$.}
\qaitem{What is $\mathbf{Prof}$?}{The bicategory of profunctors: small categories, profunctors as 1-cells (functors $\mathcal{B}^{\mathrm{op}} \times \mathcal{A} \to \mathbf{Set}$), natural transformations as 2-cells, composition via coend.}
\qaitem{What is the composition rule in $\mathbf{Prof}$?}{$(Q \circ P)(c,a) = \int^b Q(c,b) \times P(b,a)$, a coend.}
\qaitem{What is $\mathbf{Rel}$?}{The locally posetal bicategory of sets with relations as 1-cells, containments as 2-cells, and relational composition.}
\qaitem{What is $\mathbf{Bimod}$?}{The bicategory of commutative rings, $(S,R)$-bimodules as 1-cells, bimodule maps as 2-cells, and tensor product as composition.}
\qaitem{Define biequivalence.}{A pseudofunctor that is essentially surjective on objects (up to equivalence) and locally an equivalence (each hom-functor is an equivalence of categories).}
\qaitem{What is the analogue of ``fully faithful and essentially surjective'' for biequivalences?}{Locally an equivalence and essentially surjective on objects up to equivalence — the bicategorical version.}
\qaitem{What is a strict 2-functor?}{A pseudofunctor whose compositors and unitors are identities.}
\qaitem{Why can't the identity of a non-strict bicategory be strict?}{Because the associator is non-identity; strict preservation would require $F(\alpha) = \mathrm{id}$, contradicting non-strictness.}
\qaitem{In $\mathbf{Bimod}$, what is the identity 1-cell at a ring $R$?}{The bimodule $R$ itself (the regular $(R,R)$-bimodule).}
\qaitem{In $\mathbf{Span}(\mathbf{Set})$, what is the identity 1-cell at a set $A$?}{The span $A \leftarrow A \to A$ with both legs the identity.}
\qaitem{What 2-categorical structure on $\mathbf{Cat}$ is implied by the interchange law?}{Strict 2-category. The interchange law is exactly the functoriality of horizontal composition, which is what makes $\mathbf{Cat}$ strictly $\mathbf{Cat}$-enriched.}
\qaitem{Is a monoidal functor a pseudofunctor of one-object bicategories?}{Yes: a strong (= pseudo) monoidal functor is precisely a pseudofunctor of one-object bicategories; lax monoidal is lax; oplax is oplax.}
\qaitem{What is the relationship between $\mathbf{Bicat}$ as a tricategory and $\mathbf{Cat}$ as a 2-category?}{$\mathbf{Cat}$ is a sub-2-category of $\mathbf{Bicat}$: strict 2-categories with strict 2-functors and strict 2-natural transformations.}
\qaitem{In what sense does strictification leave us in no loss of generality?}{Up to biequivalence: every bicategorical statement can be checked in the strictified 2-category. The ``up to coherent iso'' caveats can be left implicit.}
\end{qa}
\end{exercisesbox}


% ═══════════════════════════════════════════════════════════════════════════
\section*{Chapter Summary}
\addcontentsline{toc}{section}{Chapter Summary}
% ═══════════════════════════════════════════════════════════════════════════

\begin{summarybox}
\textbf{Strict 2-categories.}\quad Categories enriched over $\mathbf{Cat}$:
objects, hom-categories, composition functors, with strict associativity and
unitality. $\mathbf{Cat}$ itself, locally discrete 2-categories, and
$\mathbf{Cat}/\mathcal{C}$ are examples.

\textbf{Bicategories.}\quad Generalization where associativity and unitality
hold only up to coherent invertible 2-cells (associators $\alpha$, unitors
$\lambda, \rho$) satisfying pentagon and triangle axioms. Monoidal categories
are exactly one-object bicategories.

\textbf{Pseudofunctors, lax/oplax functors.}\quad Maps of bicategories
preserving composition up to invertible 2-cells (pseudo); up to (not
necessarily invertible) 2-cells in one direction (lax) or the other (oplax).
A lax functor $\mathbf{1} \to \mathbf{Cat}$ is a monad.

\textbf{Pseudonatural transformations and modifications.}\quad Natural
transformations between pseudofunctors with invertible 2-cell components
$\eta_f$ instead of square equalities; modifications between pseudonatural
transformations. Together, bicategories form a tricategory $\mathbf{Bicat}$.

\textbf{Examples beyond $\mathbf{Cat}$.}\quad $\mathbf{Span}(\mathcal{C})$
(pullback composition), $\mathbf{Prof}$ (coend composition), $\mathbf{Rel}$
(locally posetal), $\mathbf{Bimod}$ (tensor over a ring).

\textbf{Strictification.}\quad Every bicategory is biequivalent to a strict
2-category, via the string-replacement construction. As a corollary, every
diagram of associators and unitors commutes (Mac Lane coherence for
bicategories). The strictification licenses the convention of treating
bicategories as strict in informal arguments.
\end{summarybox}


% ═══════════════════════════════════════════════════════════════════════════
\section*{Formal Restatement}
\addcontentsline{toc}{section}{Formal Restatement}
% ═══════════════════════════════════════════════════════════════════════════

\begin{formalrestatement}
\textbf{Bicategory data.}\quad Objects; hom-categories $\mathcal{B}(a,b)$;
composition functors $\circ_{a,b,c}$; identity 1-cells $1_a$; invertible
2-cells $\alpha_{h,g,f}, \lambda_f, \rho_f$. Pentagon and triangle axioms.

\medskip
\textbf{Pseudofunctor.}\quad Object map; hom-functors $F_{a,b}$;
invertible 2-cells $\phi_{g,f} : F(g) \circ F(f) \cong F(g \circ f)$ and
$\phi_a : 1_{Fa} \cong F(1_a)$; coherence with $\alpha, \lambda, \rho$.
Lax: drop invertibility. Oplax: reverse direction.

\medskip
\textbf{Pseudonatural transformation $\eta : F \Rightarrow G$.}\quad 1-cells
$\eta_a : Fa \to Ga$ and invertible 2-cells $\eta_f : Gf \circ \eta_a \cong
\eta_b \circ Ff$ satisfying compatibility with $\phi^F, \phi^G$.

\medskip
\textbf{Modification $m : \eta \Rrightarrow \eta'$.}\quad 2-cells $m_a :
\eta_a \Rightarrow \eta'_a$ compatible with $\eta_f, \eta'_f$.

\medskip
\textbf{Biequivalence.}\quad Pseudofunctor that is essentially surjective on
objects (up to equivalence) and locally an equivalence.

\medskip
\textbf{Strictification (Theorem~\ref{thm:strictification}).}\quad For every
bicategory $\mathcal{B}$ there is a strict 2-category $\mathcal{B}^{\mathrm{st}}$
and a biequivalence $\widehat{(-)} : \mathcal{B}^{\mathrm{st}} \to \mathcal{B}$.

\medskip
\textbf{Mac Lane coherence (Corollary~\ref{cor:bicat-coherence}).}\quad In any
bicategory, every diagram of associators and unitors commutes.

\medskip
\noindent\textit{Chapter~26 introduces \emph{ends and coends}, the foundational
universal constructions for working with profunctors, weighted limits, and
enriched Yoneda. Coends will be the calculus that powers most of the remaining
chapters of Volume III.}
\end{formalrestatement}
